% called by main.tex
%
\chapter{Estado de la cuestión}
\label{ch::capitulo2}

Este capítulo revisa los fundamentos y los trabajos más relevantes sobre los que
se asienta el presente Trabajo Fin de Máster. Se parte del papel del aprendizaje
profundo en el diagnóstico por imagen, se describen los conjuntos de datos y las
arquitecturas de referencia para la clasificación de patologías torácicas, y se
abordan los tres ejes centrales del trabajo: la clasificación multietiqueta en
presencia de desbalanceo, la gestión de la incertidumbre en el etiquetado y la
explicabilidad. El capítulo concluye situando la aportación del proyecto frente a
la literatura existente.


\section{El aprendizaje profundo en el diagnóstico por imagen}
\label{sec::ec_dl_imagen}

Antes de la consolidación del aprendizaje profundo, la clasificación automática
de imágenes médicas se abordaba mediante el paradigma clásico del aprendizaje
automático: un experto diseñaba manualmente un conjunto de características
(\textit{handcrafted features}) —descriptores de textura, forma o bordes— que
posteriormente alimentaban a un clasificador como las máquinas de vectores de
soporte o los bosques aleatorios. Este enfoque adolecía de una limitación
fundamental: la calidad del resultado dependía de la pericia con que se diseñaran
las características, un proceso laborioso, específico de cada problema y que rara
vez capturaba la complejidad de las imágenes médicas~\cite{litjensSurveyDeepLearning2017}.

El punto de inflexión llegó en 2012, cuando una red neuronal convolucional
(CNN) ganó con amplio margen la competición
ImageNet~\cite{krizhevskyImageNetClassification2012, dengImageNet2009}. Las CNN
sustituyen el diseño manual de características por su \textbf{aprendizaje
automático y jerárquico} directamente a partir de los datos: las primeras capas
detectan bordes y texturas elementales, y las profundas combinan esa información
en conceptos cada vez más abstractos. Esta capacidad para aprender
representaciones de extremo a extremo, unida a su invarianza ante traslaciones,
las hizo idóneas para el reconocimiento de imágenes y se trasladó con rapidez al
dominio biomédico~\cite{caiReviewApplicationDeep2020}.

La versatilidad de las CNN ha quedado patente en dominios muy diversos, desde el
diagnóstico de enfermedades de plantas~\cite{sharmaPlantDiseaseDiagnosis2022} o
la clasificación de imágenes minerales~\cite{liuEfficientImageSegmentation2021}
hasta la categorización de resonancias magnéticas
cerebrales~\cite{krishnapriyaPretrainedDeepLearning2023}, lo que evidencia la
generalidad del enfoque. En el ámbito médico, múltiples estudios comparativos
confirman que el aprendizaje profundo supera de forma consistente a los
algoritmos clásicos cuando se dispone de volúmenes de datos
suficientes~\cite{wangComparativeAnalysisImage2021, sharmaDeepLearningModels2022,
bharadiyaConvolutionalNeuralNetworks2023}.

Precisamente la dependencia de grandes volúmenes de datos constituye el principal
escollo en medicina, donde los conjuntos etiquetados son costosos y escasos. La
respuesta habitual es el \textbf{aprendizaje por transferencia} (\textit{transfer
learning}): reutilizar un modelo preentrenado sobre un gran repositorio de
imágenes naturales —como ImageNet— y adaptarlo a la tarea médica. Esta estrategia
puede aplicarse congelando el modelo preentrenado y entrenando solo la capa de
salida (\textit{feature extraction}) o reajustando también sus pesos
(\textit{fine-tuning}); en ambos casos acelera la convergencia y mejora la
generalización con menos datos~\cite{plestedDeepTransferLearning2026}. Se ha
convertido en una práctica estándar en el diagnóstico por imagen, incluida la
radiografía de tórax~\cite{jenaArtificialIntelligenceBased2021, kermanyIdentifyingMedical2018}.
No obstante, el aprendizaje profundo arrastra dos limitaciones que condicionan su
aplicación clínica y que vertebran este trabajo: su sensibilidad a las
particularidades de los datos médicos —desbalanceo, multietiqueta e
incertidumbre— y su naturaleza opaca, abordadas en las secciones siguientes.


\section{Conjuntos de datos de radiografía de tórax}
\label{sec::ec_datasets}

El rendimiento del aprendizaje profundo supervisado es directamente proporcional
a la cantidad y la calidad de los datos etiquetados disponibles. La disponibilidad
de grandes repositorios públicos de radiografías de tórax anotadas ha sido, por
tanto, el principal motor del progreso en este campo. El reto, sin embargo, es
mayúsculo: anotar manualmente cientos de miles de imágenes por parte de
radiólogos resulta inviable en coste y tiempo, lo que ha obligado a recurrir a la
extracción automática de etiquetas a partir de los informes radiológicos mediante
procesamiento del lenguaje natural (NLP), con la consiguiente introducción de
ruido.

El primer repositorio de gran escala fue \textbf{ChestX-ray14} —inicialmente
publicado como ChestX-ray8—, del \textit{National Institutes of Health} (NIH),
con más de cien mil imágenes frontales etiquetadas automáticamente para catorce
patologías~\cite{wangChestXray8Hospital2017}. Pese a su impacto, estudios
posteriores señalaron la limitada fiabilidad de sus etiquetas, derivada del
proceso automático. Le siguieron repositorios que mejoraron el tamaño o la
calidad de la anotación. \textbf{CheXpert}~\cite{irvinCheXpertLargeChest2019}, de
la Universidad de Stanford, aporta 224.316 imágenes y catorce observaciones, e
introduce como rasgo distintivo el \textbf{tratamiento explícito de la
incertidumbre} en el etiquetado. \textbf{MIMIC-CXR}~\cite{johnsonMIMICCXR2019},
del Instituto Tecnológico de Massachusetts, es aún mayor (377.110 imágenes) e
incorpora los informes de texto libre completos. \textbf{PadChest}~\cite{bustosPadChest2020},
de origen español, ofrece más de 160.000 imágenes con una taxonomía jerárquica de
174 hallazgos y la particularidad de que una cuarta parte fue anotada manualmente
por médicos.

La calidad de la anotación es, de hecho, el principal eje diferenciador entre
estos conjuntos, y admite distintos niveles. En un extremo se sitúa el etiquetado
\textbf{automático} por NLP, escalable pero ruidoso; para mejorarlo se han
desarrollado etiquetadores más sofisticados como
\textbf{CheXbert}~\cite{smitCheXbertCombining2020}, que sustituye las reglas por
un modelo de lenguaje basado en BERT combinando anotaciones automáticas y de
expertos. En el extremo opuesto se encuentra la anotación \textbf{directa por
radiólogos}, como en \textbf{VinDr-CXR}~\cite{nguyenVinDrCXR2022}, cuyas imágenes
fueron etiquetadas y localizadas mediante recuadros por radiólogos
experimentados. Esta distinción da lugar a los conceptos de referencia
\textit{silver} y \textit{gold standard} que se detallarán en el
Capítulo~\ref{sec::capitulo3}.

La elección del conjunto de datos implica, pues, un compromiso entre escala,
calidad de anotación y riqueza de las patologías. En este trabajo se ha
seleccionado CheXpert por reunir un equilibrio favorable: un tamaño que favorece
el aprendizaje de representaciones robustas, su consolidación como banco de
pruebas de referencia en la comunidad —lo que facilita la comparación con la
literatura— y, de forma determinante, la incorporación explícita del problema de
las etiquetas inciertas, que constituye uno de los ejes metodológicos del
proyecto.


\section{Arquitecturas para la clasificación de patologías torácicas}
\label{sec::ec_arquitecturas}

La arquitectura de la red es uno de los factores que más influyen en el
rendimiento, y su evolución ha seguido una trayectoria marcada por la búsqueda de
mayor profundidad sin degradar el entrenamiento. Tras el éxito de
AlexNet~\cite{krizhevskyImageNetClassification2012}, las redes
\textbf{VGG}~\cite{simonyanVeryDeepConvolutional2015} demostraron que apilar
muchas capas de filtros pequeños ($3\times3$) mejoraba el rendimiento, aunque a
costa de un elevado número de parámetros. Sin embargo, al aumentar la profundidad
surgía el problema de la \textit{desaparición del gradiente}, que dificultaba el
entrenamiento de redes muy profundas. La solución llegó con las redes residuales
\textbf{ResNet}~\cite{heDeepResidualLearning2016}, cuyas conexiones de atajo
(\textit{skip connections}) permiten que el gradiente fluya sin atenuarse,
posibilitando redes de cientos de capas. Las redes densamente conectadas
\textbf{DenseNet}~\cite{huangDenselyConnectedConvolutional2017} llevaron esta idea
más lejos: cada capa recibe como entrada las salidas de todas las anteriores, lo
que favorece la reutilización de características y la eficiencia de parámetros.
Esta arquitectura resultó especialmente adecuada para la radiografía de tórax;
el trabajo seminal \textbf{CheXNet}~\cite{rajpurkarCheXNetRadiologistLevel2017}
empleó una DenseNet-121 para alcanzar un rendimiento comparable al de los
radiólogos en la detección de neumonía sobre ChestX-ray14.

Más recientemente, el paradigma de los \textit{transformers} —originado en el
procesamiento del lenguaje natural— ha irrumpido en visión por computador. A
diferencia de las CNN, que capturan relaciones locales, los \textit{Vision
Transformers} modelan dependencias globales mediante mecanismos de
\textbf{atención}. Su aplicación directa, no obstante, exige grandes volúmenes de
datos y un elevado coste computacional. El \textbf{Swin
Transformer}~\cite{liuSwinTransformer2021} mitiga este problema con una atención
jerárquica restringida a ventanas, más eficiente para imágenes de alta
resolución. En paralelo, \textbf{ConvNeXt}~\cite{liuConvNet2020s2022} demuestra
que las CNN puras, modernizadas con decisiones de diseño propias de los
transformers, pueden igualar su rendimiento. Ambas tendencias convergen en las
\textbf{arquitecturas híbridas CNN-Transformer}, que combinan la extracción local
de las convoluciones con la captura de dependencias globales de la atención, y que
han mostrado resultados prometedores en clasificación de patologías
torácicas~\cite{singhComputerAidedDiagnosis2024, oommenConvAttentionViT2024,
fuExplainableHybridTransformer2025}.

Esta diversidad plantea una pregunta clave: ¿qué arquitectura es preferible para
la radiografía de tórax? La respuesta no es unívoca y depende del conjunto de
datos, la tarea y los recursos. Los estudios comparativos sistemáticos no
identifican un ganador absoluto: una evaluación de múltiples arquitecturas CNN
sobre radiografía de tórax situó los valores de área bajo la curva ROC (AUROC)
entre 0,83 y 0,89 según la patología y el
modelo~\cite{bressemComparingDeepLearning2020}, diferencias a menudo dentro del
margen de variabilidad experimental. A esta indefinición se añade la cuestión de
la robustez frente a perturbaciones de la entrada, crítica para la fiabilidad
clínica y objeto de estudio activo~\cite{liuComprehensiveStudyRobustness2025}.
Esta ausencia de consenso justifica precisamente el enfoque comparativo adoptado
en este trabajo, que evalúa cinco familias representativas (DenseNet, ResNet, VGG,
ConvNeXt y Swin) bajo un protocolo homogéneo.


\section{Clasificación multietiqueta y desbalanceo de clases}
\label{sec::ec_multietiqueta}

La formulación del problema condiciona toda la cadena de modelado. En la
clasificación \textit{multiclase}, cada muestra pertenece a una única categoría
entre varias mutuamente excluyentes, y se resuelve con una capa de salida
\textit{softmax} que reparte la probabilidad entre clases. El diagnóstico
radiológico, en cambio, es intrínsecamente \textbf{multietiqueta}: una misma
radiografía puede presentar simultáneamente varias patologías no
excluyentes~\cite{zhangReviewMultiLabel2014}. Por ello se aborda como un conjunto
de problemas de clasificación binaria independientes —uno por patología—,
empleando una activación \textit{sigmoide} en cada salida y una función de
pérdida de entropía cruzada binaria, que no impone exclusividad entre clases.

El principal obstáculo en este escenario es el fuerte \textbf{desbalanceo de
clases}, consecuencia directa de la distinta prevalencia epidemiológica de las
patologías: algunas aparecen en una pequeña fracción de los estudios. Sin
corrección, el modelo tiende a minimizar la pérdida prediciendo siempre la clase
negativa, lo que produce una elevada exactitud aparente pero una capacidad nula
para detectar las patologías minoritarias~\cite{lizDeepLearningMultilabel2023},
que con frecuencia son las de mayor relevancia clínica.

La literatura ofrece dos grandes familias de soluciones. Las técnicas a
\textbf{nivel de datos} reequilibran el conjunto antes del entrenamiento, mediante
sobremuestreo de las clases minoritarias, submuestreo de las mayoritarias o
aumento de datos. Las técnicas a \textbf{nivel de algoritmo} modifican el proceso
de aprendizaje: destacan la ponderación de la entropía cruzada binaria según la
prevalencia de cada clase~\cite{usharubyBinaryCrossEntropy2020}, la \textit{focal
loss} —que reduce el peso de los ejemplos fáciles—, y las funciones que optimizan
directamente el AUROC, más robustas ante el desbalanceo que las pérdidas
convencionales~\cite{yuanLargeScaleRobustDeepAUC2021}. El problema se agrava en
escenarios de \textbf{cola larga} (\textit{long-tailed}), donde numerosos
hallazgos raros coexisten con unas pocas clases muy frecuentes, lo que ha
motivado líneas de investigación e iniciativas de evaluación
específicas~\cite{holsteCXRLT2024}.

Entre estas alternativas, este trabajo adopta la \textbf{entropía cruzada binaria
ponderada} como mecanismo principal para contrarrestar el desbalanceo, por su
sencillez, su ausencia de hiperparámetros adicionales y su integración directa en
la formulación multietiqueta, frente a la mayor complejidad de las técnicas de
remuestreo o de las pérdidas alternativas. La justificación detallada de esta
elección se desarrolla en el Capítulo~\ref{sec::capitulo3}.


\section{Gestión de la incertidumbre en el etiquetado}
\label{sec::ec_incertidumbre}

La extracción automática de etiquetas mediante NLP genera, para ciertas
patologías, anotaciones \textit{inciertas}: hallazgos descritos como posibles o
no concluyentes que no pueden tratarse directamente como positivos ni como
negativos. El trabajo original de CheXpert~\cite{irvinCheXpertLargeChest2019}
formaliza este problema y propone varias estrategias para el manejo de la
etiqueta de incertidumbre (codificada con el valor $-1$), entre ellas ignorar
las muestras inciertas (\textit{U-Ignore}), asignarlas a la clase negativa
(\textit{U-Zeros}) o a la positiva (\textit{U-Ones}), o estimarlas mediante
autoentrenamiento (\textit{U-SelfTrained}) o un planteamiento multiclase
(\textit{U-MultiClass}).

Más allá del tratamiento aislado de cada etiqueta, otros trabajos explotan las
\textit{dependencias jerárquicas} entre patologías para producir predicciones
clínicamente coherentes, integrando la estructura de relaciones entre hallazgos
junto con el manejo explícito de la
incertidumbre~\cite{phamInterpretingChestXrays2021}. La comparación sistemática
de estas estrategias constituye uno de los objetivos de diseño del presente
trabajo.


\section{Inteligencia artificial explicable en imagen médica}
\label{sec::ec_xai}

La opacidad de las redes neuronales profundas, que operan como sistemas de
``caja negra'', constituye una barrera fundamental para su adopción clínica,
donde la trazabilidad de las decisiones es un requisito
ineludible~\cite{lizlopezDeepLearningInnovations2025, lizEnsemblesConvolutional2021}.
La inteligencia artificial explicable (XAI) responde a esta necesidad mediante
técnicas que hacen interpretable el razonamiento del modelo. En clasificación de
imágenes, los métodos más extendidos son los basados en mapas de activación: el
\textit{Class Activation Mapping} (CAM)~\cite{zhouLearningDeepFeatures2016} y, de
forma destacada, \textbf{Grad-CAM}~\cite{selvarajuGradCAMVisualExplanations2017},
que genera mapas de calor a partir de los gradientes que fluyen hacia la última
capa convolucional, sin requerir modificaciones en la arquitectura. Su extensión
\textbf{Grad-CAM++}~\cite{chattopadhyayGradCAMpp2018} mejora la localización
cuando coexisten múltiples instancias de una misma clase, lo que resulta
pertinente en la radiografía de tórax. Existen revisiones sistemáticas que
recopilan y categorizan el conjunto de técnicas de XAI aplicadas al diagnóstico
por imagen~\cite{patricioExplainableDeepLearning2023}.

No obstante, la fiabilidad de estos métodos no debe presuponerse. Estudios
recientes con radiólogos han mostrado que algunos mapas de saliencia pueden ser
sensibles a perturbaciones sutiles de la entrada, generando explicaciones poco
robustas o no necesariamente alineadas con la evidencia
clínica~\cite{zhangRevisitingTrustworthiness2024}. Esta cautela motiva que, en
este trabajo, las explicaciones generadas mediante Grad-CAM se interpreten como
una ayuda a la verificación por parte del especialista y no como una garantía de
corrección del modelo.


\section{Despliegue de modelos en entornos clínicos}
\label{sec::ec_despliegue}

La traslación de un modelo de investigación a la práctica clínica plantea retos
que van más allá del rendimiento predictivo. La integración en el flujo de
trabajo radiológico, la latencia de la inferencia, la monitorización continua y
la conformidad con la normativa son factores críticos para una implantación
real~\cite{korfiatisImplementingArtificial2025}. La incorporación de técnicas de
explicabilidad y de cuantificación de la incertidumbre se considera un elemento
clave para favorecer la confianza del facultativo y la integración efectiva en
el entorno asistencial.

En esta línea, la tendencia más reciente combina la información de la imagen con
datos clínicos adicionales mediante esquemas de fusión multimodal, con el
objetivo de refinar las predicciones y aproximarlas al razonamiento
diagnóstico~\cite{liReviewDeepLearning2024, yangEnhancingChestXray2025}. El
presente trabajo aborda el despliegue mediante una plataforma web que integra la
inferencia del modelo y la visualización de las explicaciones, descrita en el
Capítulo~\ref{sec::capitulo5}.


\section{Limitaciones de la literatura actual}
\label{sec::ec_limitaciones}

Pese a los avances descritos, la revisión pone de manifiesto varias limitaciones
y retos abiertos. En primer lugar, existe una amplia variedad de arquitecturas
aplicables a la radiografía de tórax —CNN clásicas, transformers y modelos
híbridos— sin que ninguna se imponga de forma concluyente; las comparativas
publicadas emplean además protocolos y métricas heterogéneos, lo que dificulta
la reproducibilidad y la comparación
directa~\cite{bressemComparingDeepLearning2020, liuComprehensiveStudyRobustness2025}.

En segundo lugar, el desbalanceo de clases y, en particular, la distribución de
cola larga de las patologías poco frecuentes siguen penalizando el rendimiento
sobre las clases minoritarias, que son a menudo las de mayor relevancia
clínica~\cite{lizDeepLearningMultilabel2023, holsteCXRLT2024}. A ello se suma la
incertidumbre inherente al etiquetado automático por NLP: las distintas
estrategias para tratar las etiquetas inciertas no presentan un ganador claro y
su efecto depende del conjunto de datos y de la
patología~\cite{irvinCheXpertLargeChest2019, phamInterpretingChestXrays2021}.

En tercer lugar, la explicabilidad, imprescindible para la adopción clínica,
no debe asumirse como fiable por defecto: estudios recientes muestran que los
mapas de saliencia pueden ser poco robustos y no siempre coinciden con la
evidencia clínica~\cite{zhangRevisitingTrustworthiness2024}. Por último, persiste
una brecha entre la investigación y la práctica: son escasos los trabajos que
abordan el despliegue \textit{de extremo a extremo} en un entorno usable,
integrando inferencia y explicabilidad, y que consideran los sesgos demográficos
de los conjuntos de datos —de origen institucional único— que limitan la
generalización del modelo~\cite{korfiatisImplementingArtificial2025,
bustosPadChest2020}.


\section{Posicionamiento del trabajo}
\label{sec::ec_sintesis}

Este Trabajo Fin de Máster se posiciona en la intersección de las limitaciones
anteriores: propone una evaluación comparativa y reproducible de arquitecturas
para la clasificación multietiqueta sobre CheXpert, con tratamiento explícito del
desbalanceo de clases y de la incertidumbre del etiquetado, e integra la
explicabilidad mediante Grad-CAM en una plataforma web desplegable que aproxima
el modelo a un uso real. La metodología y los recursos empleados para ello se
detallan en el Capítulo~\ref{sec::capitulo3}.
